\section*{Chapitre 5 - Opérations sur les nombres décimaux}
\addcontentsline{toc}{section}{Chapitre 5 - Opérations sur les nombres décimaux}

\subsection*{I. Vocabulaire}
\addcontentsline{toc}{subsection}{I. Vocabulaire}

\begin{Définition*}[c]
\leavevmode\vspace{-\baselineskip}
\begin{itemize}[label=\textcolor{Couleur6}{$\bullet$}]
\item Le résultat d'une addition s'appelle une \textbf{somme}. Le résultat d'une soustraction s'appelle une \textbf{différence}. Le résultat d'une multiplication s'appelle un \textbf{produit}.
\item Les nombres que l'on ajoute ou que l'on soustrait s'appellent des \textbf{termes}.
\item Les nombres que l'on multiplie s'appellent des \textbf{facteurs}.
\end{itemize}
\end{Définition*}

\begin{Exemples}
\begin{itemize}[label=\textcolor{Couleur8}{$\bullet$}]
\item $12 + 67 + 15 = 94$ est une somme dont les trois termes sont $12$ ; $67$ et $15$.
\item $2 \times 34 = 68$ est un produit dont les deux facteurs sont $2$ et $34$.
\end{itemize}
\end{Exemples}

\subsection*{II. Priorités opératoires}
\addcontentsline{toc}{subsection}{II. Priorités opératoires}
\subsubsection*{a. Calculs sans parenthèse}
\addcontentsline{toc}{subsubsection}{a. Calculs sans parenthèse}

\begin{minipage}[c]{0.48\textwidth}
\begin{Règle}

Dans une expression \textbf{sans parenthèse} comportant
uniquement des additions et des soustractions, on effectue les calculs \textbf{de gauche à droite}.
\end{Règle}

\begin{Exemple} 
\vspace{-0.5cm}
\begin{align*}
A & ={\color{Couleur1}20-3}+7-2\\
A & ={\color{Couleur1}17+7}-2\\
A & ={\color{Couleur1}24-2}\\
A & =22
\end{align*}

\end{Exemple}
\end{minipage} \hfill
\begin{minipage}[c]{0.48\textwidth}


\begin{Règle}
Dans une expression \textbf{sans parenthèse} comportant
uniquement des multiplications et des divisions, on effectue les calculs \textbf{de gauche à droite}.
\end{Règle}

\begin{Exemple} 
\vspace{-0.5cm}
\begin{align*}
B & ={\color{Couleur1}24:3}\times2:4\\
B & ={\color{Couleur1}8\times2}:4\\
B & ={\color{Couleur1}16:4}\\
B & =4
\end{align*}

\end{Exemple}

\end{minipage} 


\begin{Règle}
Dans une expression \textbf{sans parenthèse}, on
effectue \textbf{d'abord les multiplications et les divisions}, et
seulement après, les additions et les soustractions.
\end{Règle}


\begin{Exemples}
\begin{align*}
C & =12-{\color{Couleur1}5\times2} & D & ={\color{Couleur1}4\times5}+{\color{Couleur1}6:2}\\
C & =12-10 & D & =20+3\\
C & =2 & D & =23
\end{align*}
\end{Exemples}

\begin{Remarque}
\vspace{-0.5cm}
\begin{enumerate}
\item Dans l'expression $C$, comme le dernier calcul est une soustraction, alors $C$ est la \textbf{différence de 12 et du produit de 5 par 2.}
\item Dans l'expression $D$, comme le dernier calcul est une addition, alors $D$ est la \textbf{somme du produit de 4 par 5 et du quotient de 6 par 2.}
\end{enumerate}
\end{Remarque}


 
\subsubsection*{b. Calculs avec parenthèses}
\addcontentsline{toc}{subsubsection}{b. Calculs avec parenthèses}

\begin{Règle}
Dans une expression \textbf{avec parenthèses}, on
effectue \textbf{d'abord les calculs entre parenthèses}.
\end{Règle}

\begin{Exemples}
{\normalsize \begin{align*}
E & =6\times{\color{Couleur1}(5+3)} & F & ={\color{Couleur1}(22-2)}:4 & G & ={\color{Couleur1}(6+4)}\times{\color{Couleur1}(7-5)}\\
E & =6\times8 & F & =20:4 & G & =10\times2\\
E & =48 & F & =5 & G & =20
\end{align*}}
\end{Exemples}

\begin{Règle}
Dans une expression \textbf{avec parenthèses} \textbf{doubles},
on effectue \textbf{en premier les calculs situés dans les parenthèses les plus à l'intérieur}.
\end{Règle}

\begin{Exemple} 
{\normalsize
\begin{align*}
H & =2\times\left(8-{\color{Couleur1}\left(7-4\right)}\right)\\
H & =2\times{\color{Couleur1}\left(8-3\right)}\\
H & =2\times5\\
H & =10
\end{align*}}
\end{Exemple}


\subsection*{III. Poser les calculs}
\addcontentsline{toc}{subsection}{III. Poser les calculs}

\begin{Méthode}[c]
Pour poser une addition ou une soustraction avec des nombres décimaux, il faut aligner les chiffres des unités et rajouter des 0 pour avoir le même nombre de décimales. \\
\begin{center}
\Addition[Solution,CouleurCadre=Couleur7!10,CouleurSolution=Couleur6]{18.9}{1.255}
\hspace{2cm}
\Soustraction[Solution,CouleurCadre=Couleur7!10,CouleurSolution=Couleur6]{18.5}{12.97}
\end{center}

\end{Méthode}

\begin{Exemples}
\begin{center}
\Addition[Solution,CouleurCadre=white,CouleurSolution=Couleur6]{30.2}{2.24}
\hspace{2cm}
\Soustraction[Solution,CouleurCadre=white,CouleurSolution=Couleur6]{31.6}{11.83}
\end{center}
\end{Exemples}

\begin{Méthode}[c]
Pour poser une multiplication avec des nombres décimaux, on pose la multiplication sans écrire la ou les virgules. Une fois le calcul fini, on rajoute la virgule en comptant le nombre de décimales présentes dans l'énoncé. \\
\begin{center}
\Multiplication[Solution,CouleurCadre=Couleur7!10,CouleurSolution=Couleur6]{18.9}{1.75}
\end{center}
Ici, il  y avait 3 décimales dans l'énoncé, donc on a mis la virgule entre 33 et 075 pour qu'il  y ait bien 3 chiffres après la virgule dans le résultat.
\end{Méthode}

\begin{Exemple}
\begin{center}
\Multiplication[Solution,CouleurCadre=white,CouleurSolution=Couleur6]{6.56}{21.4}
\end{center}
\end{Exemple}


\newpage